\documentclass[a4paper,11pt]{article}

\usepackage{inputenc}
\usepackage[T1]{fontenc}
%\usepackage[frenchb]{babel}
\usepackage{fancyhdr,fancybox} % pour personnaliser les en-têtes
\usepackage{lastpage,setspace}
\usepackage{amsfonts,amssymb,amsmath,amsthm,mathrsfs}
\usepackage{mathdots}
\usepackage{relsize,exscale,bbold}
\usepackage{paralist}
\usepackage{xspace,multicol,diagbox,array}
\usepackage{xcolor}
\usepackage{variations}
\usepackage{xypic}
\usepackage{eurosym,stmaryrd}
\usepackage{graphicx}
\usepackage[np]{numprint}
\usepackage{hyperref} 
\usepackage{tikz}
\usepackage{colortbl}
\usepackage{multirow}
\usepackage{MnSymbol,wasysym}
\usepackage[top=1.5cm,bottom=1.5cm,right=1.2cm,left=1.5cm]{geometry}
\usetikzlibrary{calc, arrows, plotmarks, babel,decorations.pathreplacing}
\setstretch{1.25}
%\usepackage{lipsum} %\usepackage{enumitem} %\setlist[enumerate]{itemsep=1mm} bug avec enumerate



\newtheorem{thm}{Théorème}
\newtheorem{rmq}{Remarque}
\newtheorem{prop}{Propriété}
\newtheorem{cor}{Corollaire}
\newtheorem{lem}{Lemme}
\newtheorem{prop-def}{Propriété-définition}

\theoremstyle{definition}

\newtheorem{defi}{Définition}
\newtheorem{ex}{Exemple}
\newtheorem*{rap}{Rappel}
\newtheorem{cex}{Contre-exemple}
\newtheorem{exo}{Exercice} % \large {\fontfamily{ptm}\selectfont EXERCICE}
\newtheorem{nota}{Notation}
\newtheorem{ax}{Axiome}
\newtheorem{appl}{Application}
\newtheorem{csq}{Conséquence}
\def\di{\displaystyle}



\renewcommand{\thesection}{\Roman{section}}\renewcommand{\thesubsection}{\arabic{subsection} }\renewcommand{\thesubsubsection}{\alph{subsubsection} }


\newcommand{\bas}{~\backslash}\newcommand{\ba}{\backslash}
\newcommand{\C}{\mathbb{C}}\newcommand{\K}{\mathbb{K}}\newcommand{\R}{\mathbb{R}}\newcommand{\Q}{\mathbb{Q}}\newcommand{\Z}{\mathbb{Z}}\newcommand{\N}{\mathbb{N}}\newcommand{\V}{\overrightarrow}\newcommand{\Cs}{\mathscr{C}}\newcommand{\Ps}{\mathscr{P}}\newcommand{\Rs}{\mathscr{R}}\newcommand{\Gs}{\mathscr{G}}\newcommand{\Ds}{\mathscr{D}}\newcommand{\happy}{\huge\smiley}\newcommand{\sad}{\huge\frownie}\newcommand{\danger}{\begin{tikzpicture}[x=1.5pt,y=1.5pt,rotate=-14.2]
	\definecolor{myred}{rgb}{1,0.215686,0}
	\draw[line width=0.1pt,fill=myred] (13.074200,4.937500)--(5.085940,14.085900)..controls (5.085940,14.085900) and (4.070310,15.429700)..(3.636720,13.773400)
	..controls (3.203130,12.113300) and (0.917969,2.382810)..(0.917969,2.382810)
	..controls (0.917969,2.382810) and (0.621094,0.992188)..(2.097660,1.359380)
	..controls (3.574220,1.726560) and (12.468800,3.984380)..(12.468800,3.984380)
	..controls (12.468800,3.984380) and (13.437500,4.132810)..(13.074200,4.937500)
	--cycle;
	\draw[line width=0.1pt,fill=white] (11.078100,5.511720)--(5.406250,11.875000)..controls (5.406250,11.875000) and (4.683590,12.812500)..(4.367190,11.648400)
	..controls (4.050780,10.488300) and (2.375000,3.675780)..(2.375000,3.675780)
	..controls (2.375000,3.675780) and (2.156250,2.703130)..(3.214840,2.964840)
	..controls (4.273440,3.230470) and (10.640600,4.847660)..(10.640600,4.847660)
	..controls (10.640600,4.847660) and (11.332000,4.953130)..(11.078100,5.511720)
	--cycle;
	\fill (6.144520,8.839900)..controls (6.460940,7.558590) and (6.464840,6.457090)..(6.152340,6.378910)
	..controls (5.835930,6.300840) and (5.320300,7.277400)..(5.003900,8.554750)
	..controls (4.683590,9.835940) and (4.679690,10.941400)..(4.996090,11.019600)
	..controls (5.312490,11.097700) and (5.824210,10.121100)..(6.144520,8.839900)
	--cycle;
	\fill (7.292960,5.261780)..controls (7.382800,4.898500) and (7.128900,4.523500)..(6.730460,4.421880)
	..controls (6.328120,4.324220) and (5.929680,4.535220)..(5.835930,4.898500)
	..controls (5.746080,5.261780) and (5.999990,5.640630)..(6.402340,5.738340)
	..controls (6.804690,5.839840) and (7.203110,5.625060)..(7.292960,5.261780)
	--cycle;
	\end{tikzpicture}}\newcommand{\alors}{\Large\Rightarrow}\newcommand{\equi}{\Leftrightarrow}
\newcommand{\fonction}[5]{\begin{array}{l|rcl}
		#1: & #2 & \longrightarrow & #3 \\
		& #4 & \longmapsto & #5 \end{array}}


\definecolor{vert}{RGB}{11,160,78}
\definecolor{rouge}{RGB}{255,120,120}
\definecolor{bleu}{RGB}{15,5,107}



\pagestyle{fancy}
\lhead{Groupe IPESUP}\chead{}\rhead{Année~2022-2023}\lfoot{M. Botcazou \& M.Dupré}\cfoot{\thepage/3}\rfoot{PCSI }\renewcommand{\headrulewidth}{0.4pt}\renewcommand{\footrulewidth}{0.4pt}


\begin{document}
 %%%%BIBMATH%%%%
 
%(1) Variables aléatoires, moments https://www.bibmath.net/ressources/index.php?action=affiche&quoi=bde/proba/vagenerale&type=fexo
 
%(2)  Inégalités de Markov et de Tchebychev https://www.bibmath.net/ressources/index.php?action=affiche&quoi=bde/proba/va6-vaineg&type=fexo

\noindent\shadowbox{
	\begin{minipage}{1\linewidth}
		\centering
		\huge{\textbf{ TD 22 : Variables aléatoires et moments }}
	\end{minipage}}

\smallskip
\section*{Connaître son cours:}
\begin{itemize}[$\bullet$]
	\item Soit $X_1$ et $X_2$ deux variables indépendantes à valeurs dans $E_1$ et $E_2$ , $f_1 : E_1 \to F_1$ et $f_2 : E_2 \to F_2 $. Montrer que les variables $f_1 (X_1 )$ et $f_2 (X_2 )$ sont indépendantes.
	\item Soit $X$ une variable aléatoire définie sur un espace probabilisé $(\Omega,\mathcal{A},\mathbb{P})$ à valeurs dans $E$ et $f : E \to F$. Donner et démontrer la formule de transfert.
	\item Montrer que la somme de $n$ variables indépendantes de loi de Bernoulli de paramètre $p$ est une variable aléatoire de loi binomiale de paramètres $n$ et $p $.
	\item Soit $X$ une variable aléatoire réelle, rappeler la définition de l'espérance de $X$. Donner l'espérance d'une variable aléatoire suivant une loi binomiale de paramètre $(n,p)$. 
	\item Rappeler l'inégalité de Markov et donner la démonstration. En déduire l'inégalité de Bienaymé-Tchebychev.
	\item Montrer que la variance d'une somme de variables aléatoires indépendantes est la somme des variances de ces variables.
	\item Soient $X$ et $Y$ deux variables aléatoires réelles qui admettent un moment d'ordre 2. Démontrer la formule de Huygens.

\end{itemize}
\raggedright

\section*{Variables aléatoires discrètes et espérances:}\hfill\\%[-0.25cm]

   
\begin{minipage}{1\linewidth}\begin{minipage}[t]{0.48\linewidth}\raggedright
	
\begin{exo}\textbf{(*)}\quad\\[0.2cm] %p.
	Soit $X,Y\sim \mathcal{U}(\llbracket0;n\rrbracket)$ indépendantes.\\
	\begin{enumerate}
		\item Donner la loi de $Z_1 = X+Y$. 
		\item Donner la loi de $Z_2 = n-X$.
		\item Soit $a\in\R^*_+$, donner la loi de $Z_3 = aX$.
		\item Donner l'espérance des variables aléatoires précédentes. 
	\end{enumerate}
	
\centering\rule{1\linewidth}{0.6pt}\end{exo}



\begin{exo}\textbf{(**)}\quad\\[0.2cm]
	\begin{enumerate}
		\item Donner une variable aléatoire $X$ à valeurs dans $\N$ qui n'admet pas d'espérance finie.
		\item Soit $X$ une variable aléatoire à valeurs dans $\N$ qui admet une espérance et telle que $\mathbb{P}(X>n) \underset{\sim_{+\infty}}{=}o\left(\frac{1}{n}\right)$.
		
		 Montrer que  $$\mathbb{E}(X) = \sum_{k=0}^{+\infty}\mathbb{P}(X> k)$$
	\end{enumerate}
	%$\llbracket 0, n \rrbracket$.
	

	\centering\rule{1\linewidth}{0.6pt}\end{exo}




%%%%%%%%%%%%%%%%%%%%%%%%%%%%%%%%%%%%%%%%%%%%%%%%%%%%%%%%%%%%%%%%%%%%%%%%%%%%%%%%%%%%%%%%%%
\end{minipage}\hfill\vrule\hfill\begin{minipage}[t]{0.48\linewidth}\raggedright
%%%%%%%%%%%%%%%%%%%%%%%%%%%%%%%%%%%%%%%%%%%%%%%%%%%%%%%%%%%%%%%%%%%%%%%%%%%%%%%%%%%%%%%%%%

\begin{exo}\textbf{(*)}\quad\\[0.2cm]
Soit $X$ et $Y$ deux variables aléatoires indépendantes suivant des lois de Poisson de paramètre respectif $\lambda$ et $\mu$. Déterminer par deux méthodes différentes la loi de $Z = X+Y$ et en déduire l'espérance de $Z$. 

\centering\rule{1\linewidth}{0.6pt}\end{exo}


\begin{exo}\textbf{(**)}\quad\\[0.2cm]% Delauney p.400 MP algèbre
Soit $X$ et $Y$ deux variables aléatoires indépendantes géométriques de paramètres $p$ et $q$ à valeurs dans $]0;1[$. Calculer l'espérance de $Z= \max(X,Y)$. 

\centering\rule{1\linewidth}{0.6pt}\end{exo}

\begin{exo}\textbf{(**)}\quad\\[0.2cm]% Delauney p.400 MP algèbre
	Soit $N\in\N^*$ et $(U_n)_{n\geq1}$ une suite de variables aléatoires indépendantes et identiquement distribuées selon une loi uniforme sur $\llbracket1;N\rrbracket$. Pour tout $n\in\N^*$, on pose 
	\begin{center}
		$M_n = \max(U_1, \dots , U_n).$
	\end{center}
	Déterminer les limites de $\mathbb{E}(M_n)$ et de $\mathbb{V}(M_n)$ lorsque n tend vers $+\infty$.
	 
	\centering\rule{1\linewidth}{0.6pt}\end{exo}

\end{minipage}\end{minipage} \newpage


\begin{minipage}{1\linewidth}\begin{minipage}[t]{0.48\linewidth}\raggedright
		
		\begin{exo}\textbf{(**)}\quad\\[0.2cm] %Dupuy exo5
			On suppose qu'à la roulette d'un Casino, on obtient la couleur noire avec la probabilité $\frac{1}{2}$, la couleur rouge sinon (on ne suppose qu'il n'y a pas de 0 vert). Un joueur fortuné joue selon le protocole suivant :
			
			\begin{itemize}
				\item il mise initialement 1 brouzouf sur la couleur noire.
				\item s'il gagne, il arrête de jouer et empoche le double de sa mise.
				\item s'il perd, il double sa mise et rejoue. 
			\end{itemize}
		
			\begin{enumerate}
				\item On suppose la fortune du joueur infinie. Montrer que le jeu s'arrête presque sûrement. Déterminer l'espérance de gain du joueur.
				\item On suppose toujours la fortune du joueur infinie.	Que se passe-t-il si au lieu de doubler, il décide de tripler sa mise lorsqu'il rejoue?
				\item Le joueur n'est en fait pas si fortuné qu'il le prétend : il ne possède que $2^{n}-1$ brouzoufs ce qui l'autorise à ne pouvoir jouer que $n$ parties. Que devient son espérance de gain?
			\end{enumerate}

			
			\centering\rule{1\linewidth}{0.6pt}\end{exo}
		
		
		
		\begin{exo}\textbf{(**)}\quad\\[0.2cm]
			
			Dans une urne figurent $N$ boules numérotées de 1 à $N$ (avec $N \geq 2$ ). Dans celle-ci on opère des tirages successifs (avec remise) jusqu'à l'obtention d'une série de $k$ boules consécutives identiques $(k \geq 2)$. On admet qu'il est presque sûr que ce processus s'arrête et on note $T$ la variable aléatoire déterminant le nombre de tirages opérés à l'arrêt du processus.
			
			\begin{enumerate}
				\item Déterminer $\mathrm{P}(T=k)$ et $\mathrm{P}(T=k+1)$.
				\item Soit $n \geq 1$, établir
				
				$
				\displaystyle\mathrm{P}(T=n+k)=\frac{N-1}{N^{k}} \mathrm{P}(T>n)
				$
				\item En déduire que la variable $T$ admet une espérance et déterminer celle-ci.
				
			\end{enumerate}  
			
			\centering\rule{1\linewidth}{0.6pt}\end{exo}
		
		\begin{exo}\textbf{(**)}\quad\\[0.2cm]
			
			Soit $X$ une variable aléatoire suivant une loi géométrique de paramètre $p$. 
			
			Calculer $
			\mathrm{E}\left(\dfrac{1}{X}\right) \text {. }
			$
			
			\centering\rule{1\linewidth}{0.6pt}\end{exo}
		
	
		
		
		%%%%%%%%%%%%%%%%%%%%%%%%%%%%%%%%%%%%%%%%%%%%%%%%%%%%%%%%%%%%%%%%%%%%%%%%%%%%%%%%%%%%%%%%%%
	\end{minipage}\hfill\vrule\hfill\begin{minipage}[t]{0.48\linewidth}\raggedright
		%%%%%%%%%%%%%%%%%%%%%%%%%%%%%%%%%%%%%%%%%%%%%%%%%%%%%%%%%%%%%%%%%%%%%%%%%%%%%%%%%%%%%%%%%%
		
		\begin{exo}\textbf{(**)}\quad\\[0.2cm]
			On lance une pièce équilibrée jusqu'à ce que celle-ci ait produit au moins une fois «face et une fois «pile ».
			
			\begin{enumerate}
				\item Justifier qu'il est presque sûr que le jeu s'arrête.
				\item On note $X$ le nombre de lancers avant que le jeu cesse. Montrer que $X$ admet une espérance et déterminer celle-ci.
			\end{enumerate}
			
			
			
			\centering\rule{1\linewidth}{0.6pt}\end{exo}
		
		
		\begin{exo}\textbf{(**)}\quad\\[0.2cm]
			
			Soient $X$ et $Y$ deux variables aléatoires à valeurs dans $\mathbb{N}$.
			On suppose que la loi conjointe de $X$ et $Y$ vérifie
			$$
			\displaystyle\forall(j, k) \in \mathbb{N}^{2}, \mathrm{P}(X=j, Y=k)=a \frac{j+k}{2^{j+k}} \text { avec } a \in \mathbb{R} .
			$$
			\begin{enumerate}
				\item Déterminer la valeur de $a$ pour que la somme des probabilités soit égale à 1.
				\item Déterminer les lois marginales de $X$ et de $Y$. 
				\item Les variables $X$ et $Y$ sont elles indépendantes?
				\item Calculer $\mathrm{P}(X=Y)$.
			\end{enumerate}
			
			\centering\rule{1\linewidth}{0.6pt}\end{exo}
		
		\begin{exo}\textbf{(***)}\quad\\[0.2cm]
			
			Soit $T$ une variable aléatoire à valeurs naturelles vérifiant pour tout $n \in \mathbb{N}, \quad \mathrm{P}(T>n)>0$.
			
			On appelle taux de panne associé à $T$ la suite $\left(\theta_{n}\right)_{n \in \mathbb{N}}$ déterminée par
			$$
			\theta_{n}=\mathrm{P}(T=n \mid T \geq n)
			$$
			
			Typiquement, si $T$ est la variable aléatoire indiquant l'instant où un matériel tombe à panne, la quantité $\theta_{n}$ indique la probabilité qu'il tombe en panne à l'instant présent alors qu'il est actuellement fonctionnel.
			
			\begin{enumerate}
				\item Justifier $ \forall n \in \mathbb{N}, \theta_{n} \in[0 ; 1[$	
				\item Exprimer en fonction des termes de la suite $\left(\theta_{n}\right)_{n \in \mathbb{N}}$, la probabilité $\mathrm{P}(T \geq n)$. En déduire la divergence de la série $\sum \theta_{n}$.
			\end{enumerate}

			
			
			\centering\rule{1\linewidth}{0.6pt}\end{exo}
		
\end{minipage}\end{minipage}

\section*{Moments d'ordre 2 d'une variable aléatoire:}\hfill\\%[-0.25cm]


\begin{minipage}{1\linewidth}\begin{minipage}[t]{0.48\linewidth}\raggedright
		
		\begin{exo}\textbf{(*)}\quad\\[0.2cm] %Mansuy p.971
		Soit $X$ une variable aléatoire admettant un moment d'ordre 2. Faire une étude de la fonction $a\mapsto E\big((X-a)^2\big)$ et donner une interprétation de la variance de $X$. 
						
			\centering\rule{1\linewidth}{0.6pt}\end{exo}
		
				\begin{exo}\textbf{(*)}\quad\\[0.2cm]
			On dit qu'une variable aléatoire discrète réelle $X$ est quasi-certaine lorsqu'il existe un réel $a$ tel que $P(X=a)=1$. Soit $X$ une variable aléatoire discrète réelle telle que $X(\Omega)$. Démontrer que $X$ est quasi-certaine si et seulement si $V(X)=0$.
			
			\centering\rule{1\linewidth}{0.6pt}\end{exo}
		
		
		\begin{exo}\textbf{(**)}\quad\\[0.2cm]
			Soient $X$ et $Y$ deux variables aléatoires réelles admettant chacune une variance. On suppose $\mathbb{V}(X)>0$. Déterminer $a, b \in \mathbb{R}$ minimisant la quantité
			
			$$
			\mathrm{E}\left((Y-(a X+b))^{2}\right)
			$$
			\centering\rule{1\linewidth}{0.6pt}\end{exo}
		
		
		
		
		%%%%%%%%%%%%%%%%%%%%%%%%%%%%%%%%%%%%%%%%%%%%%%%%%%%%%%%%%%%%%%%%%%%%%%%%%%%%%%%%%%%%%%%%%%
	\end{minipage}\hfill\vrule\hfill\begin{minipage}[t]{0.48\linewidth}\raggedright
		%%%%%%%%%%%%%%%%%%%%%%%%%%%%%%%%%%%%%%%%%%%%%%%%%%%%%%%%%%%%%%%%%%%%%%%%%%%%%%%%%%%%%%%%%%
		

		
		
		\begin{exo}\textbf{(**)}\quad\\[0.2cm] %EX0 3 Dupuy
			
			Soit $X$ une variable aléatoire discrète à valeurs \\
			dans $[a ; b]$.
			\begin{enumerate}
				\item  Montrer que $X$ admet une espérance $m$ et que celle-ci est élément de $[a ; b]$. La variable $X$ admet aussi une variance $\sigma^{2}$ que l'on se propose de majorer. On introduit la variable aléatoire $Y=X-m$ et les quantités\end{enumerate}
$$
t=\sum_{y \geq 0} y P(Y=y),\  s=\sum_{y \geq 0} y^{2} \mathrm{P}(Y=y) \text { et }  \ u=\mathrm{P}(Y \geq 0) .
$$			\begin{enumerate}\setcounter{enumi}{1}
				\item Vérifier $t^{2} \leq s u$
				\item Calculer espérance et variance de $Y$.\\ 
				En déduire $t^{2} \leq\left(\sigma^{2}-s\right)(1-u)$
				\item En exploitant les deux majorations précédentes, obtenir
				$$
				t^{2} \leq \sigma^{2} / 4
				$$
				\item Conclure $$
				\sigma^{2} \leq(b-a)^{2} / 4
				$$
			\end{enumerate}
	

			\centering\rule{1\linewidth}{0.6pt}\end{exo}
		
		
\end{minipage}\end{minipage}\section*{Inégalités de Markov et de Bienaymé-Tchebychev :}\hfill\\%[-0.25cm]


\begin{minipage}{1\linewidth}\begin{minipage}[t]{0.48\linewidth}\raggedright



\begin{exo}\textbf{(*)}\quad\\[0.2cm]
	Soit $X$ une variable aléatoire réelle finie à valeurs dans $\mathbb R_+$, $f:\mathbb R_+\to\mathbb R_+^*$ une fonction croissante. Démontrer que
	
	$P(X\geq a)\leq \frac{E(f(X))}{f(a)}.$	\centering\rule{1\linewidth}{0.6pt}\end{exo} 


\begin{exo}\textbf{(**)}\quad\\[0.2cm]
	On jette 3600 fois un dé équilibré. Minorer la probabilité que le nombre d'apparitions du numéro 1 soit compris entre 480 et 720.
	
	\centering\rule{1\linewidth}{0.6pt}\end{exo}

\begin{exo}\textbf{(**)}\quad\\[0.2cm]
	Une usine fabrique des pièces dont une proportion inconnue $p$ est défectueuse, et on souhaite trouver une valeur approchée de $p$. 
	\end{exo}







%%%%%%%%%%%%%%%%%%%%%%%%%%%%%%%%%%%%%%%%%%%%%%%%%%%%%%%%%%%%%%%%%%%%%%%%%%%%%%%%%%%%%%%%%%
\end{minipage}\hfill\vrule\hfill\begin{minipage}[t]{0.48\linewidth}\raggedright
%%%%%%%%%%%%%%%%%%%%%%%%%%%%%%%%%%%%%%%%%%%%%%%%%%%%%%%%%%%%%%%%%%%%%%%%%%%%%%%%%%%%%%%%%%
On effectue un prélèvement de $n$ pièces. On suppose que le prélèvement se fait sur une population très grande, et donc qu'il peut s'apparenter à une suite de $n$  tirages indépendants avec remise. On note $X_n$ la variable aléatoire égale au nombre de pièces défectueuses et on souhaite quantifier le fait que $\dfrac{X_n}{n}$ approche $p$.
\begin{enumerate}
	\item Quelle est la loi de $X_n$? Sa moyenne? Sa variance?
	\item Démontrer que, pour tout $\varepsilon>0$,\\[0.5cm]
	\begin{center}	
	 $\displaystyle P\left(\left|\frac{X_n}n-p\right|\geq\varepsilon\right)\leq\frac 1{4n\varepsilon^2}.$
	\end{center}
	\item En déduire une condition sur $n$ pour que $X_n/n$ soit une valeur approchée de $p$ à $10^{-2}$ près avec une probabilité supérieure ou égale à $95\%$.
\end{enumerate}
	
	\centering\rule{1\linewidth}{0.6pt}




\end{minipage}\end{minipage} \newpage

\end{document}